\documentclass[11pt]{article}
% Preambulo comun para los informes de la ronda 3.
% Se incluye con % Preambulo comun para los informes de la ronda 3.
% Se incluye con % Preambulo comun para los informes de la ronda 3.
% Se incluye con \input{../comun/preambulo} desde cada subcarpeta.
\usepackage[utf8]{inputenc}
\usepackage[T1]{fontenc}
\usepackage[spanish,es-nodecimaldot]{babel}
\usepackage{lmodern}
\usepackage[a4paper,margin=2.4cm]{geometry}
\usepackage{booktabs}
\usepackage{longtable}
\usepackage{array}
\usepackage{xcolor}
\usepackage{hyperref}
\usepackage{enumitem}
\usepackage{fancyhdr}
\usepackage{titlesec}
\usepackage{graphicx}
\usepackage{amssymb}
\usepackage{microtype}

\definecolor{sosTeal}{HTML}{0F766E}
\definecolor{sosGray}{HTML}{475569}

\hypersetup{
  colorlinks=true,
  linkcolor=sosTeal,
  urlcolor=sosTeal,
  citecolor=sosTeal,
  pdftitle={SOS Clinic - Ronda 3 de pruebas},
  pdfauthor={Equipo SOS Clinic}
}

\pagestyle{fancy}
\fancyhf{}
\lhead{\textcolor{sosGray}{SOS Clinic --- Ronda 3}}
\rhead{\textcolor{sosGray}{\leftmark}}
\cfoot{\thepage}
\setlength{\headheight}{15pt}
\renewcommand{\headrulewidth}{0.4pt}

\titleformat{\section}{\Large\bfseries\color{sosTeal}}{\thesection}{0.8em}{}
\titleformat{\subsection}{\large\bfseries\color{sosGray}}{\thesubsection}{0.6em}{}

\setlist[itemize]{leftmargin=1.4em,itemsep=2pt,topsep=4pt}
\setlist[enumerate]{leftmargin=1.6em,itemsep=2pt,topsep=4pt}

\newcommand{\pantalla}[1]{\texttt{\detokenize{#1}}}
\newcommand{\ok}{\textcolor{sosTeal}{\textbf{PASS}}}
\newcommand{\fallo}{\textcolor{red!70!black}{\textbf{FAIL}}}
\newcommand{\aviso}{\textcolor{orange!80!black}{\textbf{WARN}}}
\newcommand{\omitido}{\textcolor{sosGray}{\textbf{SKIP}}}
 desde cada subcarpeta.
\usepackage[utf8]{inputenc}
\usepackage[T1]{fontenc}
\usepackage[spanish,es-nodecimaldot]{babel}
\usepackage{lmodern}
\usepackage[a4paper,margin=2.4cm]{geometry}
\usepackage{booktabs}
\usepackage{longtable}
\usepackage{array}
\usepackage{xcolor}
\usepackage{hyperref}
\usepackage{enumitem}
\usepackage{fancyhdr}
\usepackage{titlesec}
\usepackage{graphicx}
\usepackage{amssymb}
\usepackage{microtype}

\definecolor{sosTeal}{HTML}{0F766E}
\definecolor{sosGray}{HTML}{475569}

\hypersetup{
  colorlinks=true,
  linkcolor=sosTeal,
  urlcolor=sosTeal,
  citecolor=sosTeal,
  pdftitle={SOS Clinic - Ronda 3 de pruebas},
  pdfauthor={Equipo SOS Clinic}
}

\pagestyle{fancy}
\fancyhf{}
\lhead{\textcolor{sosGray}{SOS Clinic --- Ronda 3}}
\rhead{\textcolor{sosGray}{\leftmark}}
\cfoot{\thepage}
\setlength{\headheight}{15pt}
\renewcommand{\headrulewidth}{0.4pt}

\titleformat{\section}{\Large\bfseries\color{sosTeal}}{\thesection}{0.8em}{}
\titleformat{\subsection}{\large\bfseries\color{sosGray}}{\thesubsection}{0.6em}{}

\setlist[itemize]{leftmargin=1.4em,itemsep=2pt,topsep=4pt}
\setlist[enumerate]{leftmargin=1.6em,itemsep=2pt,topsep=4pt}

\newcommand{\pantalla}[1]{\texttt{\detokenize{#1}}}
\newcommand{\ok}{\textcolor{sosTeal}{\textbf{PASS}}}
\newcommand{\fallo}{\textcolor{red!70!black}{\textbf{FAIL}}}
\newcommand{\aviso}{\textcolor{orange!80!black}{\textbf{WARN}}}
\newcommand{\omitido}{\textcolor{sosGray}{\textbf{SKIP}}}
 desde cada subcarpeta.
\usepackage[utf8]{inputenc}
\usepackage[T1]{fontenc}
\usepackage[spanish,es-nodecimaldot]{babel}
\usepackage{lmodern}
\usepackage[a4paper,margin=2.4cm]{geometry}
\usepackage{booktabs}
\usepackage{longtable}
\usepackage{array}
\usepackage{xcolor}
\usepackage{hyperref}
\usepackage{enumitem}
\usepackage{fancyhdr}
\usepackage{titlesec}
\usepackage{graphicx}
\usepackage{amssymb}
\usepackage{microtype}

\definecolor{sosTeal}{HTML}{0F766E}
\definecolor{sosGray}{HTML}{475569}

\hypersetup{
  colorlinks=true,
  linkcolor=sosTeal,
  urlcolor=sosTeal,
  citecolor=sosTeal,
  pdftitle={SOS Clinic - Ronda 3 de pruebas},
  pdfauthor={Equipo SOS Clinic}
}

\pagestyle{fancy}
\fancyhf{}
\lhead{\textcolor{sosGray}{SOS Clinic --- Ronda 3}}
\rhead{\textcolor{sosGray}{\leftmark}}
\cfoot{\thepage}
\setlength{\headheight}{15pt}
\renewcommand{\headrulewidth}{0.4pt}

\titleformat{\section}{\Large\bfseries\color{sosTeal}}{\thesection}{0.8em}{}
\titleformat{\subsection}{\large\bfseries\color{sosGray}}{\thesubsection}{0.6em}{}

\setlist[itemize]{leftmargin=1.4em,itemsep=2pt,topsep=4pt}
\setlist[enumerate]{leftmargin=1.6em,itemsep=2pt,topsep=4pt}

\newcommand{\pantalla}[1]{\texttt{\detokenize{#1}}}
\newcommand{\ok}{\textcolor{sosTeal}{\textbf{PASS}}}
\newcommand{\fallo}{\textcolor{red!70!black}{\textbf{FAIL}}}
\newcommand{\aviso}{\textcolor{orange!80!black}{\textbf{WARN}}}
\newcommand{\omitido}{\textcolor{sosGray}{\textbf{SKIP}}}


\title{\textbf{DAST con OWASP ZAP}\\[0.3em]
\large Metodología --- Ronda 3}
\author{Equipo SOS Clinic}
\date{22 de octubre de 2026}

\begin{document}
\maketitle

\section{Objetivo}

Analizar dinámicamente el frontend estático en busca de problemas de
configuración, cabeceras, exposición de información y patrones pasivos de
riesgo, sin lanzar ataques activos ni afectar datos.

\section{Alcance}

\begin{itemize}
  \item Objetivo: \texttt{http://host.docker.internal:5500}, que corresponde al
  servidor del frontend en el puerto 5500 de la máquina anfitriona.
  \item 17 URLs descubiertas por el \textit{spider} de ZAP.
  \item Perfil \textbf{baseline} (pasivo): no se ejecutan ataques de inyección
  ni escaneos activos.
  \item Sin autenticación: se analiza la superficie pública (landing, login,
  recuperación de clave y páginas estáticas).
\end{itemize}

\section{Herramienta}

\begin{table}[h]
\centering
\small
\begin{tabular}{@{}ll@{}}
\toprule
Elemento & Detalle \\
\midrule
Imagen & \texttt{ghcr.io/zaproxy/zaproxy:stable} \\
Script & \texttt{zap-baseline.py} \\
Perfil & Baseline pasivo \\
Objetivo & \texttt{http://host.docker.internal:5500} \\
Fecha & 22 de octubre de 2026 \\
\bottomrule
\end{tabular}
\caption{Configuración de la ejecución de ZAP.}
\end{table}

\section{Procedimiento}

\begin{enumerate}
  \item Verificar que un contenedor alcanza el frontend local mediante
  \texttt{host.docker.internal}.
  \item Ejecutar ZAP con el perfil baseline, limitando el \textit{spider} a un
  minuto y esperando el análisis pasivo.
  \item Generar reportes HTML, JSON y Markdown en
  \pantalla{seguridad/resultados/}.
  \item Clasificar cada advertencia por impacto real y contexto de la
  aplicación.
\end{enumerate}

\begin{verbatim}
docker run --rm \
  -v <repo>/seguridad/resultados:/zap/wrk:rw \
  ghcr.io/zaproxy/zaproxy:stable zap-baseline.py \
  -t http://host.docker.internal:5500 \
  -r zap-baseline-r3.html \
  -J zap-baseline-r3.json \
  -w zap-baseline-r3.md
\end{verbatim}

\section{Criterios}

ZAP clasifica cada regla como \texttt{PASS}, \texttt{WARN} o \texttt{FAIL}.
Un \texttt{FAIL} indica una vulnerabilidad o configuración claramente insegura;
un \texttt{WARN} requiere evaluación manual. En esta ronda se revisó cada
advertencia y se documentó si era un riesgo real, una decisión intencional o un
residual conocido.

\section{Limitaciones}

\begin{itemize}
  \item El escaneo es pasivo: no prueba inyección, lógica de negocio ni
  autenticación.
  \item No se autenticó una sesión, por lo que las páginas internas no se
  analizaron con datos reales.
  \item El frontend es estático y carga datos por JavaScript; ZAP no ejecuta
  todos los flujos de la SPA.
  \item El entorno usa HTTP local; las reglas de HSTS no aplican.
  \item No sustituye el pentest manual ni una revisión de código.
\end{itemize}

\section{Reproducción}

El comando anterior deja los tres reportes. Para revisar el resultado sin
navegador, el JSON se puede resumir con las reglas y sus ocurrencias.

\end{document}
